% year: תשפ"ו
% semester: א
% moed: א
% number: 1א
% points: 10
% categories: func-analysis, ivt
% official_solution: yes
% source: exams/מבחנים/תשפו/AnyScanner_02_01_2026(1) (2).pdf

\begin{question}
כמה פתרונות ממשיים יש למשוואה הבאה?
\[ x^4 + 2x^2 = 1 \]
נמקו היטב.
\end{question}

\begin{hint}
הציבו $t = x^2$ וקבלו משוואה ריבועית ב-$t$; זכרו ש-$t = x^2$ חייב להיות אי-שלילי.
\end{hint}

\begin{solution}
\textbf{דרך 1 (הצבה).} נעביר אגפים: $x^4 + 2x^2 - 1 = 0$. נציב $t = x^2 \ge 0$ ונקבל
\[ t^2 + 2t - 1 = 0 \quad\Longrightarrow\quad t_{1,2} = \frac{-2 \pm \sqrt{4 + 4}}{2} = -1 \pm \sqrt{2}. \]
הפתרון $t = -1 - \sqrt{2} < 0$ נפסל, כי $x^2 \ge 0$ לכל $x$ ממשי, ולכן אין לו פתרון ממשי.
הפתרון $t = -1 + \sqrt{2} > 0$ (כי $\sqrt 2 > 1$) נותן
\[ x^2 = \sqrt{2} - 1 \quad\Longrightarrow\quad x_{1,2} = \pm\sqrt{\sqrt{2} - 1}, \]
שני מספרים ממשיים שונים.

\textbf{דרך 2 (חקירה).} נגדיר $f(x) = x^4 + 2x^2 - 1$, פונקציה רציפה וגזירה בכל $\R$. מתקיים
\[ f'(x) = 4x^3 + 4x = 4x(x^2 + 1), \]
ומכיוון ש-$x^2 + 1 > 0$, הסימן של $f'$ הוא הסימן של $x$: $f' < 0$ ב-$(-\infty, 0)$ ו-$f' > 0$ ב-$(0, \infty)$.
לכן $f$ יורדת ממש ב-$(-\infty, 0]$ ועולה ממש ב-$[0, \infty)$, ובכל אחד מהקטעים האלה יש לה לכל היותר אפס אחד.
כעת $f(0) = -1 < 0$ ו-$f(1) = f(-1) = 2 > 0$. לפי משפט ערך הביניים יש ל-$f$ אפס ב-$(-1, 0)$ ואפס ב-$(0, 1)$, ולפי המונוטוניות אלה האפסים היחידים.

\textbf{תשובה:} למשוואה יש בדיוק \textbf{שני} פתרונות ממשיים, $x = \pm\sqrt{\sqrt 2 - 1}$.
\end{solution}
